\section{A polynomial estimator for the cell functional}
\label{sec:angular}

We construct an estimator of the signed entropy of an ordinary
projection \(q\) on \(\Omega\), evaluated on \(W\).
In this section, \(w\ge0\) is in \(L^2(\Omega)\), with mass
\(p_0\), and \(e=\|w-q\|_2\). Fix an integer \(J\ge1\) and set
\(\alpha=3/4\), \(\epsilon=\alpha^J/2\).
Write \(K=\lfloor c_K\log n\rfloor\), with \(c_K>0\)
sufficiently small. We use two observation models.
\begin{enumerate}
\item For a Poisson correction cloud with intensity \(m w(x)\dd x\)
  on \(\Omega\), with \(m\asymp n\), assume \(p_0\le A\), where
  \(A=C_A\log n/n\). The cloud is the restriction of a parent
  sample of mean \(m\). Use \(a_*=A\) and \(r_*=A/K^2\).
\item For \(N\) observations from the probability density \(w\),
  with \(C_*\log n\le N\le n\), use an interval pilot and
  an independent correction fraction, each of size comparable
  to \(N\). Here \(p_0=1\), \(a_*=1\),
  \(t=\log n/N\) and \(r_*=t/K^2\).
\end{enumerate}
For fixed \(J\) and any \(\eta>0\), we will estimate
\(\int_W [\psi(q)-q\log a_*]\) with absolute bias at most
\begin{equation}\label{eq:angular-bias}
 C\left[e+r_*\left(1+(p_0/r_*)^\epsilon
                      \one_{\{\inf_W|q|<p_0/2\}}\right)\right].
\end{equation}
The second moment bounds will be \(C n^\eta A^2\log^2 n\) in
the Poisson model and \(C n^\eta\log^2 n\) in the fixed-count
model. In the latter model the bounds condition on the interval
pilot on an event with probability at least \(1-n^{-M}\)
for any prescribed \(M\). We describe the event and the statistic
on arbitrary pilot histories.

\subsection{Regional scores and angular intervals}

For a Borel set \(S\subset\Omega\) with \(v_S=|S|>0\),
we need a polynomial reproducing kernel on \(S\) with bounded
observation scores. On the polynomial space of dimension \(D\)
use the norm \(v_S^{-1}\int_S|P|\), which is definite by
Lemma~\ref{lem:sublevel}. Maximize the absolute determinant of
\(D\) vectors in its compact unit ball. The maximizing vectors
\(\varphi_i\) have norm one, since the maximum is positive.
Replacing one by another vector in the ball shows that its
coordinate functional has norm at most one.
Extend these functionals by Hahn--Banach to \(L^1\) of normalized
volume on \(S\). Duality gives representatives
\(\sigma_i\) bounded by one, with
\(v_S^{-1}\int_S\sigma_i\varphi_j=\delta_{ij}\).
Use finite Borel versions, clipped to \([-1,1]\) on \(S\)
and extended by zero off \(S\). Changing a Lebesgue null set
preserves these pairings and the score means in both models.

Let
\[
 B_S=\sum_i|\varphi_i|,\qquad
 \ell_S=B_S/v_S,\qquad
 K_S(x,y)=v_S^{-1}\one_S(y)\sum_i\varphi_i(x)\sigma_i(y).
\]
Since the constant one has all coordinates of magnitude at most
one, \(B_S\ge1\). Also \(\int_S B_S=Dv_S\).
The kernel reproduces every polynomial in the space and obeys
\(|K_S(x,y)|\le\ell_S(x)\one_S(y)\). Define
\(q_S=\int K_S(\cdot,y)w(y)\dd y\), \(c_S=q_S/\ell_S\)
and \(p_S=\int_S w\). Subtracting the reproduction of \(q\)
inside the integral and using Cauchy--Schwarz gives
\begin{equation}\label{eq:regional}
 |c_S|\le p_S,\qquad \int_S\ell_S=D,\qquad
 \|q_S-q\|_{1,S}\le D\sqrt{v_S}\,e .
\end{equation}
The actual scores \(g_{S,0}=\one_S\) and
\(g_{S,i}=\one_S(1+\sigma_i)/2\) belong to \([0,1]\)
everywhere. Their means against \(w\) are in \([0,p_S]\).
Using \(2\int w g_{S,i}-p_S\) for the signed coordinate
makes \(c_S(x)\) a linear form in the means with total absolute
coefficient bounded by three. For the ordinary root
\(S_0=\Omega\) use \(q_0=q\) and a constant envelope
\(\ell_0\ge\max(1,\|K_\Omega\|_\infty)\). The ordinary basis
scores are affine rescalings of the bounded basis polynomials
to \([0,1]\), together with the mass indicator. Their linear
form for \(c_0=q/\ell_0\) has bounded coefficient sum, and
\(|c_0|\le p_0\). The first kernel uses \(\Omega\) and its
first cut acts on \(X_0=W\). Each later kernel uses the
remaining part of that evaluation set.

In the Poisson model each score mean has interval \([0,A]\).
In the other model an empirical mean \(\widehat u\) gives
\[
 [l,l+B]=[(\widehat u-\Delta)_+,\widehat u+\Delta],
 \qquad \Delta=L_1(\sqrt{t\widehat u}+t).
\]
These intervals always have width at least \(L_1t\).
For a true mean \(u\), Bernstein's inequality uses
\(0\le g\le1\) and \(\Var(g)\le u\) to give
\(|\widehat u-u|\le C(\sqrt{tu}+t)\) at any fixed
inverse-polynomial failure probability. This implies
\(\sqrt{t\widehat u}+t\asymp\sqrt{tu}+t\), by separating
\(u\le C't\) from \(u>C't\) with \(C'\) sufficiently large.
Thus, with \(L_1\) large enough,
\begin{equation}\label{eq:interval}
 u\in[l,l+B],\qquad B\asymp\sqrt{tu}+t
\end{equation}
on that event. The fixed dictionaries of scores used below
have polynomial size in \(n\), so this bound holds simultaneously
with probability at least \(1-n^{-M}\) by a union bound.
Every interval lies in a fixed bounded nonnegative range,
including on failing histories.

Write \(u=l+B(1+\cos\theta)/2\) for one interval of width \(B\).
The derivative satisfies \(|u'(\theta)|^2=(u-l)(l+B-u)\).
In either model, valid intervals give the increment bound
\begin{equation}\label{eq:increment}
 |u(\theta+z)-u(\theta)|
 \le C\left(\sqrt{r_*p_S}\,K|z|+r_*(K|z|)^2\right)
\end{equation}
for a score supported on \(S\). In the Poisson model
\(B=A\), \(l=0\), and the first derivative has magnitude at most
\(\sqrt{Au}\) while the second derivative is \(O(A)\).
For the fixed-count model and \(u\ge t\), use
\(B|z|\le C\sqrt{tu}|z|\).
For \(u<t\) use \(B\le Ct\), first derivative bounded by
\(\sqrt{B(u-l)}\) and second derivative \(O(B)\).
The same increment estimate with a different constant applies
to \(c_S(x)\), uniformly in \(x\) and with a norm of the angle
vector in place of \(|z|\). Set \(D_S=r_*+\sqrt{r_*p_S}\).

\subsection{Polynomial selection and approximation}

On the angle circle, use a nonnegative even trigonometric
polynomial density of degree at most \(K\). For any fixed
integer \(b\ge1\) and large \(K\), we can take the normalized
\(2b\)-th power of the modulus of
\(\sum_{k=0}^{\lfloor K/b\rfloor}e^{ikz}\).
Its bound on \([-\pi,\pi]\) is
\(C_b K(1+K|z|)^{-2b}\). The geometric sum is bounded by
\(C\min(K,|z|^{-1})\) and has magnitude at least \(c_b K\)
on an interval of length \(c_b/K\), which yield the bound.
To define path probabilities at the true means, shift the
angles independently by this density, transform back to the
score intervals, and round on a nonnegative mesh of spacing at
most \(n^{-3}\). The statistic uses the polynomial path
probabilities in its coefficients, as shown below.
This size is \(O(r_*)\) throughout our count ranges.

For any rounded vector its probability, as a function of true
angles, is convolution of the product shift density with the
indicator of a bin. The bin is invariant under sign change in
each angle, so the probability is an even trigonometric
polynomial in each angle, hence a polynomial in the normalized
means \((u-l)/B\), of degree \(O(K)\). The sum of the
coefficient norms over all bins is at most \(C^K\).
Indeed the Fourier coefficients of the shift density, using
normalized Haar measure, have modulus at most one.
Each Fourier coefficient of a bin indicator is bounded by
its Haar measure, and those bin measures sum to one.
The number of Fourier indices is polynomial in \(K\).
The monomial coefficient norm of
\(\cos(k\theta)=T_k(2(u-l)/B-1)\) is at most \(C^k\)
by the Chebyshev recurrence. Multiplying over the fixed
number of angles gives the summed bound. Here and below the
coefficient norm is the sum of absolute monomial coefficients
in the normalized means.

Use a separate signed filter for approximation. In each score
angle its Fourier multiplier is \(\chi(k/K)\), where
\(\chi\) is even and smooth, supported in \([-1,1]\) and
equal to one on \([-1/2,1/2]\). Use the product filter in all
angles. Its kernel is the periodization of a rescaled Schwartz
function on Euclidean space, by Fourier inversion, and it
reproduces every trigonometric polynomial whose coordinate
degrees are at most \(K/2\).
For \(G(z)=\psi(z)-z\log r_*\), applying this filter to
\(G(c_S(x))\) as a function of angles produces a polynomial
\(L_{S,x}\) in the normalized means, of degree \(O(K)\),
with the following bound at the true angles on valid intervals
for any fixed \(b'>0\):
\begin{equation}\label{eq:filter}
 |L_{S,x}-G(c_S(x))|\le
        C_{b'} D_S(1+|c_S(x)|/D_S)^{-b'}.
\end{equation}
To prove the bound, write \(c=c_S(x)\) at the true means
and \(v\) for the increment due to the angle shift.
With \(z/K\) denoting the Euclidean shift before
periodization, \eqref{eq:increment} bounds \(|v|\) by
\(CD_S(|z|+|z|^2)\).
If \(|c|\le D_S\), subtract the exactly reproduced linear term
\(z'\log(D_S/r_*)\) in the function argument \(z'\), and
rescale the remaining \(z'\log|z'/D_S|\) by \(D_S\).
The Schwartz bound gives error \(O(D_S)\). If \(|c|>D_S\),
subtract the Taylor polynomial in \(v\) at \(c\) up to some
fixed order \(a-1\), with \(a\ge2\). Its powers of the increment
have finite angle degree at most \(a-1\) and are reproduced
for large \(K\). On \(|v|\le |c|/2\) the error is bounded by
\(C_a|v|^a|c|^{1-a}\). On the complement,
\(|z|+|z|^2\ge c'|c|/D_S\). After cancellation of the linear
part containing \(\log(|c|/r_*)\), the remaining difference,
including the Taylor terms, is bounded by \(|c|\) times a fixed
polynomial in \(1+|v|/|c|\). This follows by scaling
\(\psi(c+v)\) by \(|c|\) and bounding
\(|u\log|u||\le C(1+|u|^2)\). Choosing \(a\) and the
Schwartz decay order large proves \eqref{eq:filter}.
The approximation is even in each angle, which justifies
the polynomial representation.

On all rectangle arguments, \(|c_S(x)|\le C a_*\) in
either model. The Fourier coefficients of \(G(c_S(x))\)
are then at most \(C a_*(1+|\log(r_*/a_*)|)\) in
absolute value. As for the selection probabilities, conversion
from angle modes gives coefficient norm at most
\(C^K a_*(1+|\log(r_*/a_*)|)\), uniformly in \(x\).

\subsection{The recursive partition and its bias}

At stage \(j<J\), use the positive selection at \(S_j\) to
produce the rounded mass \(\widehat p_j\) and the rounded
projection \(a_j\), for which \(a_j/\ell_j\) uses the linear
form for \(c_j=c_{S_j}\). Set
\[
 B_j=L_2 r_*(1+\widehat p_j/r_*)^\alpha,\qquad
 T_j=\{x\in X_j:|a_j(x)|>\ell_j(x)B_j\},\qquad
 S_{j+1}=X_j\setminus T_j .
\]
Use \(X_j=S_j\) for \(j\ge1\), take \(T_J=S_J\),
and let null regions have zero contribution.
All possible cuts from any rounded vector in the fixed
bounded mesh range give a finite dictionary before observing
the interval-pilot sample. Its cardinality is polynomial in \(n\), since
the dimension and \(J\) are fixed. Define the regional kernels
there in advance, independently of the pilot and correction
observations. The dictionary can depend on \(n,N\).
On \(T_j\), the identity
\(\ell_j G(q/\ell_j)+q\log(\ell_j r_*/a_*)=\psi(q)-q\log a_*\)
decomposes the target using the envelope \(\ell_j\) of the
current kernel. We approximate the \(G\) term at \(c_j=q_{S_j}/\ell_j\)
and account for its difference from \(q/\ell_j\) below.
The resulting polynomial contribution on \(T_j\) is
\begin{equation}\label{eq:piece}
 \int_{T_j}\ell_j(x)L_{S_j,x}\dd x+
 \int_{T_j}q(x)\log(\ell_j(x)r_*/a_*)\dd x .
\end{equation}
The last term is linear in the ordinary global score means.
For the weighted envelope, \(\ell_S\ge1\) and
\[
 \int_S\log\ell_S
   = \int_S\log B_S+v_S\log(1/v_S)
   \le C v_S+v_S\log(1/v_S)\le C.
\]
This also gives a uniform bound for its coefficient norm after
including \(\log(r_*/a_*)\). The coefficient functions of the observed endpoints and the
angles are Borel. In particular the bin indicator uses the
specified rounding map, while \(G\) and the filter are continuous.
The subsequent operations are Fourier integrals over the angle
torus and spatial integrals of Borel integrands on predetermined
regions. These define Borel coefficients for the polynomial in
the score means once the real pilot has been observed.

To check the bias, condition the artificial selections on the
preceding rounded vectors, and put \(u_j=p_{S_j}/r_*\).
At each stage the probability that the scaled projection or mass
moves by more than \(C r_*(1+u_j)^\alpha\) is at most
\(C_M(1+u_j)^{-M}\), for any prescribed fixed \(M\).
In fact, \eqref{eq:increment} shows that for large \(u_j\)
such a move requires \(K|z|\ge c(1+u_j)^{(\alpha-1/2)/2}\)
in some angle, and the positive kernel has tails of arbitrarily
large fixed order by choice of \(b\). Rounding contributes
\(O(r_*)\). Bounded \(u_j\) are included by increasing \(C_M\).
Call the complementary event a good selection.
On it \(1+\widehat p_j/r_*\asymp1+u_j\) by \(\alpha<1\).
Choosing \(L_2\) sufficiently large, a good selection gives
\[
 |c_j|\ge c r_*(1+u_j)^\alpha\ \text{on }T_j,\qquad
 |c_j|\le C r_*(1+u_j)^\alpha\ \text{on } S_{j+1}.
\]
Consequently \eqref{eq:filter}, with \(b'\) sufficiently large,
and \(\int_{X_j}\ell_j\le C\) bound the mean filter error at
each nonterminal stage by \(C r_*\). The bad selection term
uses \(C D_{S_j}(1+u_j)^{-M}\le C r_*\).

Also on a good selection,
\[
 p_{S_{j+1}}\le C r_*(1+u_j)^\alpha+Ce.
\]
This follows by integrating \(|q_{S_j}|\) over \(S_{j+1}\),
using the bound on \(|c_j|\) and \(\int_{X_j}\ell_j\le C\),
then using \eqref{eq:regional} and
\(\|w-q\|_{1,S_j}\le\sqrt{v_{S_j}}e\).
On other selections \(p_{S_{j+1}}\le p_{S_j}\).
For any \(0<b_0\le1\) these bounds give the conditional estimate
\[
 \E[(1+u_{j+1})^{b_0}\mid\text{preceding vectors}]
 \le C\bigl((1+u_j)^{\alpha b_0}+(1+e/r_*)^{b_0}\bigr).
\]
The terminal filter error is bounded by \(C\E D_{S_J}\),
where \(D_{S_J}=r_*(1+\sqrt{u_J})\).
Iterate the conditional bound starting with \(b_0=1/2\),
reducing the exponent by \(\alpha\) at each prior stage.
After \(J\) steps the exponent on \(1+u_0\) is \(\epsilon\).
All residual powers have exponent at most one, so each
contributes \(O(r_*+e)\) after multiplication by \(r_*\).
The terminal error is at most \(C[r_*(1+u_0^\epsilon)+e]\). Terminated paths
can be assigned zero mass and zero contributions.
If \(\inf_W|q|\ge p_0/2\), for sufficiently large \(u_0\)
the first good selection empties \(S_1\) since
\(r_*(1+u_0)^\alpha=o(p_0)\). On the bad first selection
use \(D_{S_J}\le D_{S_0}\), together with its small
probability. Bounded \(u_0\) give the same \(O(r_*)\) bound
directly. Thus the terminal power loss is confined to
\(\inf_W|q|<p_0/2\).

It remains to replace \(q_{S_j}\) by \(q\) in the nonlinear
term of \eqref{eq:piece}. On an interval
\([-C p_0,C p_0]\), replace \(G\) by zero on
\([-r_*,r_*]\) and leave it equal to \(G\) outside.
The resulting function is continuous with Lipschitz constant
\(C(1+\log(1+u_0))\) and differs from \(G\) by at most \(C r_*\).
Both \(q_{S_j}/\ell_j\) and \(q/\ell_j\) lie in that interval.
After multiplication by \(\ell_j\) and integration on \(T_j\),
\eqref{eq:regional} bounds the error of the replacement by
\[
 C r_*+ C(1+\log(1+u_0))\sqrt{|S_j|}\,e .
\]
At the ordinary root \(q_0=q\). On a good first selection,
\(|q|\le C r_*(1+u_0)^\alpha\) on \(S_1\); for large \(u_0\)
apply Lemma~\ref{lem:sublevel} to \(q/p_0\), of mean one, to
bound \(|S_1|\le C(1+u_0)^{-(1-\alpha)\beta}\).
On a bad first selection the probability bound suffices, since
\(|S_j|\le1\). The expected replacement error is therefore
\(O(r_*+e)\). Bounded \(u_0\), including \(p_0=0\),
need only the volume bound. The rescaling identity used to define \eqref{eq:piece} now
applies on every piece of the partition. Summing the mean
errors proves \eqref{eq:angular-bias}.

\subsection{Estimating the path polynomial}

Average the contributions \eqref{eq:piece} exactly with their
polynomial path probabilities, multiplying by the probability
of each selection up to and including the one defining \(T_j\).
No selections are sampled in the implemented estimate.
At each node the outgoing coefficient norms sum to at most
\(C^K\); submultiplicativity and the fixed depth give total
degree at most \(C_J K+1\) and total coefficient norm at most
\begin{equation}\label{eq:path-C}
 C_J^K a_*(1+|\log(r_*/a_*)|).
\end{equation}
This bounds the sum of coefficients over the entire path
expansion, including products of scores on different regions.

To estimate a monomial, use the ordered distinct products
of the actual normalized observations \(Y=(g-l)/B\).
In the fixed-count correction fraction of size \(m_c\), divide
the degree-\(k\) ordered sum by \((m_c)_k=m_c(m_c-1)
\cdots(m_c-k+1)\). It is unbiased, conditional on the real
pilot. On valid intervals \(|\E Y|\le1\) and
\(\E Y^2\le1+C/t\), since \(\Var(g)\le u\) and
\(B\asymp\sqrt{tu}+t\).
In the Poisson model normalize \(Y=g/A\).
If \(p_0\le A<1\), a single parent observation may be represented
by the law on \(\Omega\) with density \(w\) and extra mass
\(1-p_0\) at one point outside \(\Omega\), where all scores
are zero. Consequently \(\E Y\le1\) and \(\E Y^2\le1/A\),
using \(g^2\le g\) pointwise. The Poisson ordered sum is divided
by \(m^k\), where \(m\) is the deterministic parent mean.
It can use just the restricted cloud since tuples with an
outside score vanish. The expectation is the product of the
parent score means, by the factorial moments of \(\Pois(m)\).
Equivalently, the local count law is \(\Pois(m p_0)\), and for
\(p_0>0\) its marks have conditional density \(w/p_0\), yielding
the same expectation. The degree-zero term is one in either
model, including with an empty cloud.

Here is a second moment bound that also applies to mixed scores.
Let \(B_2\) bound their individual second moments. Two ordered
\(k\)-tuples sharing \(i\) observations have pair contribution
at most \(B_2^i\) by Cauchy--Schwarz on each common observation
and the bound on unshared means. In the fixed-count group,
the pairing count after division by the squared denominator
is at most \(\binom{k}{i}(2k/m_c)^i\) for \(k\le m_c/2\).
For the second tuple, choose the positions of the common points,
assign distinct points from the first tuple, and bound the
remaining choices by \((m_c-i)_{k-i}\). Division by
\((m_c)_k\) yields the bound. The Poisson count instead gives
the factor \(m^{2k-i}/m^{2k}\) for the distinct points,
so the bound holds with \(2k/m\) in place of \(2k/m_c\).
Thus the second moment is at most \((1+2k B_2/m_c)^k\)
or its Poisson analogue. In our applications
\(mA\asymp\log n\) and \(m_c t\asymp\log n\).
Choose \(c_K\) so \(C_J K+1\le m_c/2\) at the
smallest allowed \(N\). The bounds on the monomials and
Minkowski with \eqref{eq:path-C} give second moments
\(C_J^K A^2\log^2 n\) and \(C_J^K\log^2 n\), respectively,
with \(C_J\) increased if necessary.
For any prescribed \(\eta>0\), further decrease \(c_K\)
to bound the exponential factor by \(C n^\eta\).
Our estimator is finite on arbitrary allowed actual pilot
histories, because the interval widths are positive and the
computed paths use only the finite set of Borel scores.
